\subsection{Kalshi event contracts}

Kalshi is a US exchange for event contracts, regulated by the Commodity Futures Trading Commission
as a designated contract market. Each binary contract pays \$1 if its event occurs and \$0
otherwise, and is fully collateralised. Trading takes place in a central limit order book. A bid
for NO at price $q$ is economically an offer of YES at $1-q$, so each market has a single YES book.
During our sample, prices were quoted on a one-cent grid from 1\cents to 99\cents.

Kalshi charges a taker fee per order of $0.07\times C\times P(1-P)$, rounded up, for $C$ contracts at
price $P$, with series-specific multipliers. Some series also carry a maker fee with coefficient
$0.0175$. Both fees vanish at the edges of the price range. Together with the one-cent tick, this
shapes trading near 0 and 1 and matters for interpreting small departures from fair pricing
(\cref{sec:discussion}).

Contracts are organised as \emph{series} (a recurring question), \emph{events} (one instance, such
as one game or one city-day) and \emph{markets} (one binary leg). Several of our designs exploit
this structure. A game has a winner leg and ladders of spread and total legs. A city-day of
temperature has a ladder of threshold or bracket legs. A question such as ``Will a model be
released before $T$'' is often listed with several deadlines.

\subsection{Scope}

We study \emph{non-financial} contracts in Kalshi's seven categories:
\begin{itemize}
  \item Sports;
  \item Climate and Weather;
  \item Entertainment;
  \item Mentions (whether a speaker says a phrase during an event);
  \item Science and Technology;
  \item Politics;
  \item Elections.
\end{itemize}
We exclude crypto, equity-index, interest-rate and economic-release contracts. Their prices are
tied to traded underlyings with their own dynamics, whereas the events we study have no traded
underlying, so the information structure of the contract itself is the object of interest.

\subsection{Data}

All data come from Kalshi's public market-data interface. We use no account-level information.

\paragraph{Market metadata.} 1.30 million binary markets in the seven categories that closed between
15 August and 31 October 2026. For each market we have:
\begin{itemize}
  \item its series and event;
  \item the rules text;
  \item open and close times;
  \item the scheduled expected expiration;
  \item the settlement result;
  \item the fee schedule.
\end{itemize}

\paragraph{Hourly quotes.} Hourly candles with the closing best bid and ask. The base panel has
60,231 markets with the highest volume in each category, from 13 February to 31 October 2026. Two
studies supplement it:
\begin{itemize}
  \item the ladder study fetches complete ladders, 1.76 million additional hourly candles;
  \item the deadline study fetches the full history of deadline contracts whose deadlines fall
  between 1 January and 2 November 2026.
\end{itemize}

\paragraph{Trades.} 7.3 million trade prints on 14,771 markets between 30 July and 31 October 2026,
each with size, price and aggressor side.

\paragraph{Order-book recordings.} We recorded the full WebSocket feed: order-book snapshots, deltas
and trades. Receive timestamps were taken on a server in the exchange's cloud region. The
recordings cover two evenings with many live games: 7 November 17:00 to 8 November 06:00 UTC, and
8 November 17:00 to 9 November 06:00 UTC. They give event-time order books for in-game
game-winner markets (NBA, NCAAF, NFL, NHL).

\subsection{Timing conventions and look-ahead}

Three timestamps in the metadata could bias the analysis if used naively.
\begin{itemize}
  \item \emph{Close time.} The recorded close time of a market that resolves early reveals the event.
  We never use it to define a horizon.
  \item \emph{Scheduled settlement.} The scheduled \emph{expected expiration} is announced in advance,
  but it marks settlement, not the end of information. It lies 11--14 hours after the last weather
  observation and days after many speeches.
  \item \emph{Deadlines.} The deadline of a ``by $T$'' contract is parsed from its rules text.
\end{itemize}
Each study therefore defines time to expiry $\tauexp$ from information available at the decision
time:
\begin{itemize}
  \item the single-contract study uses a per-series scheduled end of information, estimated on the
  estimation sample only;
  \item the ladder study uses leads before the scheduled settlement;
  \item the deadline study uses the parsed deadline.
\end{itemize}
Quotes are always taken from the last candle (or book state) strictly before the decision time.

\subsection{Sample split}

Every parameter, functional form and threshold is chosen on the \emph{estimation sample}, data
before 1 October 2026. Results are then reported on the \emph{holdout}, October 2026, without
re-estimation. The order-book study covers two November evenings and is purely descriptive.
